\begin{ouverture}[espace]{17}{Pyramides et cônes}{Partie 4 · Géométrie et espace}
{Reconnaître · Perspective · Patron · Volume}
\objectif{Reconnaître une pyramide, un cône de révolution et les autres solides usuels.}
\objectif{Relier perspective cavalière et patron d’une pyramide ou d’un cône.}
\objectif{Calculer les volumes d’une pyramide et d’un cône.}
\prerequis{Aires usuelles, volumes du pavé, du prisme et du cylindre, Pythagore.}
\end{ouverture}
\pagestyle{otMain}\essentiel\fontsize{11}{15}\selectfont
\phantomsection\label{cours-17-1}
\essSection{1}{Reconnaître et représenter}
\begin{definition}Une pyramide a une base polygonale et des faces latérales triangulaires qui se rejoignent en un sommet. Un cône de révolution est obtenu en faisant tourner un triangle rectangle autour d’un côté de l’angle droit : sa base est un disque, sa hauteur joint le sommet au centre de la base et une génératrice joint le sommet au cercle de base.\end{definition}
\begin{savoirfaire}[Distinguer les solides]Un prisme droit ou un cylindre a deux bases superposables et parallèles. Une pyramide ou un cône a une seule base et un sommet. La base d’une pyramide peut être un triangle, un carré ou un autre polygone.\end{savoirfaire}
\label{fig-cours-17-1-cours-solides}\figureGeo{cours-solides}

\begin{savoirfaire}[Perspective et patron]La perspective cavalière dessine les arêtes cachées en pointillés ; elle ne conserve pas toutes les longueurs ni tous les angles. Un patron est une figure plane qui reconstitue le solide par pliage : les côtés assemblés doivent avoir la même longueur. Pour un cône, l’arc du secteur latéral a la longueur du cercle de base.\end{savoirfaire}
\clearpage
\phantomsection\label{cours-17-2}
\essSection{2}{Des volumes liés à une base}
\begin{propriete}Pour une pyramide, $V=\frac13\mathcal B\times h$, où $\mathcal B$ est l’aire de la base et $h$ la hauteur, perpendiculaire au plan de base. Pour un cône de rayon $r$, $V=\frac13\pi r^2h$. Il faut le contenu de trois pyramides pour remplir un prisme de même base et de même hauteur ; de même pour un cône et un cylindre.\end{propriete}
\begin{exemple}Une pyramide de base carrée de côté 6 cm et de hauteur 8 cm a pour volume $\frac13\times36\times8=96$ cm$^3$. Un cône de rayon 3 cm et de hauteur 5 cm a pour volume $\frac13\pi\times9\times5=15\pi$ cm$^3$.\end{exemple}
\begin{savoirfaire}[Choisir les bonnes données]Calculer l’aire de la base avant le volume. La hauteur est perpendiculaire à la base : ce n’est ni une arête inclinée ni une génératrice. Convertir les longueurs dans une même unité avant de les multiplier : une aire en cm$^2$ multipliée par une hauteur en cm donne un volume en cm$^3$.\end{savoirfaire}
\begin{attention}La perspective peut écraser une hauteur ou un angle droit : se fier aux données et aux codages, sans mesurer le dessin. Pour un cône, diviser le diamètre par 2 avant d’utiliser le rayon dans $\pi r^2$.\end{attention}
\begin{exemple}[Repère culturel]La pyramide du Louvre, à Paris, est une pyramide de verre à base carrée ; les pyramides de Gizeh, en Égypte, ont aussi une base carrée. Le défi D17 étudie un modèle de la première.\end{exemple}

\clearpage\rubrique{Découverte}
\begin{blocLivre}{ex-17-1}\begin{exercice}[1]{Nommer les éléments}
Une pyramide SABCD a pour base le carré ABCD. Nommer son sommet, sa base et ses faces latérales.
\label{fig-ex-17-1}\figureGeo{sol-01}
\end{exercice}\end{blocLivre}
\begin{blocLivre}{ex-17-2}\begin{exercice}[1]{Une pyramide à base carrée}
Une pyramide a une base carrée de côté 6 cm et une hauteur de 8 cm. Calculer son volume.
\label{fig-ex-17-2}\figureGeo{sol-02}
\end{exercice}\end{blocLivre}
\begin{blocLivre}{ex-17-3}\begin{exercice}[1]{Un cône}
Un cône a pour rayon 3 cm et pour hauteur 5 cm. Donner son volume exact, puis arrondi au dixième de cm$^3$.
\label{fig-ex-17-3}\figureGeo{sol-03}
\end{exercice}\end{blocLivre}
\begin{blocLivre}{ex-17-4}\begin{exercice}[1]{Du diamètre au rayon}
Un cône a une base de diamètre 8 cm et une hauteur de 9 cm. Calculer son volume exact.
\label{fig-ex-17-4}\figureGeo{sol-04}
\end{exercice}\end{blocLivre}
\begin{blocLivre}{ex-17-5}\begin{exercice}[1]{Reconnaître des solides}
Nommer chacun des quatre solides représentés et préciser le nombre et la forme de ses bases.
\label{fig-ex-17-5}\figureGeo{sol-05}
\end{exercice}\end{blocLivre}

\rubrique{Entraînement}
\begin{blocLivre}{ex-17-6}\begin{exercice}[2]{Un cône par rotation}
Le triangle ABC est rectangle en A, avec $AB=4$ cm et $AC=3$ cm. On le fait tourner autour de la droite $(AB)$, puis autour de la droite $(AC)$. Décrire les deux cônes obtenus et comparer leurs volumes.
\label{fig-ex-17-6}\figureGeo{sol-06}
\end{exercice}\end{blocLivre}
\begin{blocLivre}{ex-17-7}\begin{exercice}[2]{Base triangulaire}
Une pyramide a pour base un triangle dont un côté mesure 8 cm et la hauteur associée 3 cm. La hauteur de la pyramide vaut 10 cm. Calculer son volume.
\label{fig-ex-17-7}\figureGeo{sol-07}
\end{exercice}\end{blocLivre}
\begin{blocLivre}{ex-17-8}\begin{exercice}[2]{Retrouver une hauteur}
Une pyramide a un volume de 150 cm$^3$ et une aire de base de 45 cm$^2$. Déterminer sa hauteur.
\label{fig-ex-17-8}\figureGeo{sol-08}
\end{exercice}\end{blocLivre}
\begin{blocLivre}{ex-17-9}\begin{exercice}[2]{Un patron de pyramide}
On veut construire le patron d’une pyramide à base carrée de côté 4 cm dont les quatre arêtes latérales mesurent 5 cm. Décrire puis construire ce patron.
\label{fig-ex-17-9}\figureGeo{sol-09}
\end{exercice}\end{blocLivre}
\begin{blocLivre}{ex-17-10}\begin{exercice}[2]{Un patron de cône}
Un cône a pour rayon de base 3 cm et pour génératrice 5 cm. Représenter ce cône en perspective, puis compléter son patron : un disque de base et un secteur de rayon 5 cm dont on déterminera l’angle.
\label{fig-ex-17-10}\figureGeo{sol-10}
\end{exercice}\end{blocLivre}

\rubrique{Problèmes}
\begin{blocLivre}{ex-17-11}\begin{exercice}[3]{Une génératrice calculée}
Un cône a une hauteur de 12 cm et un rayon de 5 cm. Calculer sa génératrice, puis son volume exact.
\label{fig-ex-17-11}\figureGeo{sol-11}
\end{exercice}\end{blocLivre}
\begin{blocLivre}{ex-17-12}\begin{exercice}[3]{Comparer un cylindre et un cône}
Un cône et un cylindre ont la même base de rayon 4 cm et la même hauteur de 6 cm. Donner leurs volumes et leur différence.
\label{fig-ex-17-12}\figureGeo{sol-12}
\end{exercice}\end{blocLivre}
\begin{blocLivre}{ex-17-13}\begin{exercice}[3]{Deux unités}
Un réservoir conique a un rayon de 40 cm et une hauteur de 1,2 m. Donner sa capacité en litres, au dixième.
\label{fig-ex-17-13}\figureGeo{sol-13}
\end{exercice}\end{blocLivre}
\begin{blocLivre}{ex-17-14}\begin{exercice}[3]{Un assemblage}
Un solide est formé d’un pavé droit de 8 cm sur 6 cm sur 4 cm, surmonté d’une pyramide dont la base est la face supérieure du pavé et dont la hauteur vaut 3 cm. Calculer le volume total.
\label{fig-ex-17-14}\figureGeo{sol-14}
\end{exercice}\end{blocLivre}
\begin{blocLivre}{ex-17-15}\begin{exercice}[3]{Un remplissage partiel}
Un récipient conique de rayon 6 cm et de hauteur 10 cm est rempli aux trois quarts de sa capacité. Calculer le volume d’eau exact en cm$^3$, puis en mL au dixième. Peut-on en déduire une hauteur d’eau de 7,5 cm ?
\label{fig-ex-17-15}\figureGeo{sol-15}
\end{exercice}\end{blocLivre}

\clearpage
